% Источники: exam/materials/преподаватель/2020/27-04-2020-Model_L10_3_.pdf; exam/materials/моделирование.pdf, раздел 34; GitHub HumsterProgrammer/iu9-notes/8sem/Моделирование/Моделирование_лекция-13_26-03-26.md; GitHub HumsterProgrammer/iu9-notes/8sem/Моделирование/Моделирование_лаб_30-03-26.md; /Users/egor/Desktop/сем8/Моделирование/рк/РкМод/варианты/вариант2.pdf.
\section{Коэффициенты коллигации, контингенции и ассоциации.}

\subsection*{Постановка}

Применяются для оценки связи между двумя \textbf{бинарными} (дихотомическими) признаками $A$ и $B$.
По выборке строится таблица сопряжённости $2\times 2$:

\begin{center}
\begin{tabular}{c|cc|c}
  \hline
  & $B=1$ & $B=0$ & \\
  \hline
  $A=1$ & $a$ & $b$ & $a+b$ \\
  $A=0$ & $c$ & $d$ & $c+d$ \\
  \hline
  & $a+c$ & $b+d$ & $n$ \\
  \hline
\end{tabular}
\end{center}

\subsection*{Коэффициент коллигации Юла}

\[
  Q = \frac{\sqrt{ad} - \sqrt{bc}}{\sqrt{ad} + \sqrt{bc}}.
\]

\subsection*{Коэффициент контингенции Пирсона}

\[
  \varphi = \frac{ad - bc}{\sqrt{(a+b)(c+d)(a+c)(b+d)}}.
\]

\subsection*{Коэффициент ассоциации}

\[
  A = \frac{ad - bc}{ad + bc}.
\]

\textbf{Иерархия чувствительности:} по абсолютной величине обычно $|A| \ge |Q| \ge |\varphi|$ (ассоциация наиболее чувствительна к слабым связям).

\subsection*{Алгоритм}

\begin{enumerate}
  \item Бинаризовать каждый признак (например, по медиане или по смыслу задачи).
  \item Построить таблицу сопряжённости $2\times 2$, подсчитав $a$, $b$, $c$, $d$.
  \item Вычислить все три коэффициента по формулам выше.
  \item Интерпретировать: значение $\pm1$ — максимальная связь, $0$ — связи нет.
\end{enumerate}

\subsection*{Критерий финального ответа}

В ответе должны быть построены таблицы сопряжённости и проверены соотношения из лекции: по абсолютной величине коэффициент ассоциации не меньше коэффициента коллигации, коэффициент коллигации не меньше коэффициента контингенции. Отдельно нужно учитывать, какие признаки сравниваются: количественные или качественные; при переходе от количественных признаков к бинарным требуется бинаризация.

\subsection*{Пример из варианта 2}

Имеются результаты опроса 10 человек. Обозначим ответ <<Да>> как 1, <<Нет>> как 0.

\begin{center}
\small
\begin{tabular}{c|cccccccccc}
  \hline
  Номер & 1 & 2 & 3 & 4 & 5 & 6 & 7 & 8 & 9 & 10 \\
  \hline
  $Q_1$ & Да & Нет & Да & Нет & Нет & Нет & Нет & Да & Нет & Да \\
  $Q_2$ & Да & Да & Нет & Да & Да & Нет & Да & Нет & Нет & Нет \\
  $Q_3$ & Нет & Нет & Нет & Да & Нет & Нет & Нет & Нет & Да & Нет \\
  $Q_4$ & Нет & Да & Нет & Нет & Да & Да & Нет & Нет & Нет & Да \\
  \hline
\end{tabular}
\end{center}

Для каждой пары вопросов строится таблица сопряжённости с обозначениями:
$a$ — оба ответа <<Да>>, $b$ — первый <<Да>>, второй <<Нет>>, $c$ — первый <<Нет>>, второй <<Да>>, $d$ — оба <<Нет>>.
Далее для каждой пары подставляем $a,b,c,d$ в лекционные формулы:
\[
  A=\frac{ad-bc}{ad+bc}, \qquad
  Q=\frac{\sqrt{ad}-\sqrt{bc}}{\sqrt{ad}+\sqrt{bc}}, \qquad
  \varphi=\frac{ad-bc}{\sqrt{(a+b)(c+d)(a+c)(b+d)}}.
\]
Например, для пары $Q_1,Q_2$: $a=1$, $b=3$, $c=4$, $d=2$, поэтому
\[
  A=\frac{1\cdot2-3\cdot4}{1\cdot2+3\cdot4}=-0{,}714,
\]
\[
  Q=\frac{\sqrt{1\cdot2}-\sqrt{3\cdot4}}{\sqrt{1\cdot2}+\sqrt{3\cdot4}}=-0{,}420,
  \qquad
  \varphi=\frac{1\cdot2-3\cdot4}{\sqrt{(1+3)(4+2)(1+4)(3+2)}}=-0{,}408.
\]

\begin{center}
\small
\begin{tabular}{c|cccc|ccc}
  \hline
  Пара & $a$ & $b$ & $c$ & $d$ & $A$ & $Q$ & $\varphi$ \\
  \hline
  $Q_1,Q_2$ & 1 & 3 & 4 & 2 & $-0{,}714$ & $-0{,}420$ & $-0{,}408$ \\
  $Q_1,Q_3$ & 0 & 4 & 2 & 4 & $-1{,}000$ & $-1{,}000$ & $-0{,}408$ \\
  $Q_1,Q_4$ & 1 & 3 & 3 & 3 & $-0{,}500$ & $-0{,}268$ & $-0{,}250$ \\
  $Q_2,Q_3$ & 1 & 4 & 1 & 4 & $0{,}000$ & $0{,}000$ & $0{,}000$ \\
  $Q_2,Q_4$ & 2 & 3 & 2 & 3 & $0{,}000$ & $0{,}000$ & $0{,}000$ \\
  $Q_3,Q_4$ & 0 & 2 & 4 & 4 & $-1{,}000$ & $-1{,}000$ & $-0{,}408$ \\
  \hline
\end{tabular}
\end{center}

Сильнейшая отрицательная связь по коэффициентам ассоциации и коллигации наблюдается для пар $Q_1,Q_3$ и $Q_3,Q_4$: $A=-1{,}000$, $Q=-1{,}000$, $\varphi=-0{,}408$.

\clearpage
